\section{铁磁质}
概括起来说，铁磁质有下列一些特殊的性质：
\begin{enumerate}
	\item 能产生特别强的附加磁场$B'$，使铁磁质中的$B$远大于$B_0$,其$\mu_r=B/B_0$可达几百、甚至几千。
	\item 它们的磁化强度$M$和磁感应强度$B$不再是常矢量，没有简单的正比关系，其磁导率$\mu$（以及磁化率$\chi_m$）与磁场强度$H$的函数关系比较复杂。
	\item 磁化强度随外磁场而变，它的变化落后于外磁场的变化，而且在外磁场停止作用后，铁磁质仍能保留部分磁性。
	\item 一些铁磁质存在一特定的临界温度，称为居里点（Curie point)，在温度达到居里点后，它们的磁性发生突变，转化为顺磁质，磁导率（或磁化率）和磁场强度$H$无关。例如，铁的居里点是1040$K$，镍的是631$K$，钻的是1388$K$。
\end{enumerate}